\section{Related Work}

\subsection{Mathematical Information Retrieval}

Early Math IR focused on formula-aware and keyword search, motivated by the fact that mathematical notation cannot be treated as ordinary text. Benchmarks such as NTCIR-11 Math-2, NTCIR-12 MathIR, and ARQMath formalized retrieval over mathematical corpora containing both formulas and surrounding text, including arXiv, Wikipedia, and Math Stack Exchange \cite{ntcir-11-math-2, ntcir-12-math-ir, mansouri-arqmath-3}. These settings showed that mathematical meaning depends on both symbolic structure and visual layout, motivating representations such as Symbol Layout Trees and Operator Trees with formula-specific indexing methods \cite{zanibbi2015tangentsearchengineimproved}. More recent work explores learned representations: zbMath-BERT studies embeddings for mathematical text \cite{zbmath-bert}, NaturalProofs and TheoremSearch apply neural retrieval to informal statements \cite{naturalproofs, theoremsearch2026}, and LeanSearch retrieves Mathlib declarations from informal queries via sloganized formal statements \cite{leansearch-v1, leansearch-v2}. Our work extends sloganized statement retrieval by incorporating explicit dependency structure among mathematical statements.

\subsection{Bibliometrics}

Bibliometric work uses scholarly signals to measure influence, relatedness, and field structure. Citation graphs model relationships among papers, authors, and journals; citation-count measures such as the h-index and journal impact factor summarize influence directly, while recursive prestige models such as Pinski--Narin influence, Eigenfactor, CiteRank, and \.{Z}yczkowski's weighted impact factors weight citations by graph structure, recency, or linking behavior \cite{journal-impact-factor, h-index-paper, pinski-narin, bergstrom-eigenfactor, walker-citerank, zyczkowski-citation-graph}. In mathematics, graph-based analyses have been applied to Mathlib dependency structure and proof-assistant libraries as heterogeneous networks \cite{li-network-structure-mathlib, bauer-datasets-for-formal}, while TheoremKB and AutoMathKG construct theorem-level graphs for informal sources at smaller scale \cite{mishra-theoremkb, automathkg}. Our work extends these graph-based perspectives to theorem-level dependencies across both informal and formal mathematics.

\subsection{Neural Theorem Proving \& Autoformalization}
\label{sec:ntp-autoform}
Neural theorem proving and autoformalization address complementary stages in converting informal mathematics into machine-checked proofs: autoformalization translates informal statements into formal declarations, while theorem proving generates proofs for those declarations. Recent progress on benchmarks such as miniF2F and PutnamBench has increased the importance of retrieval and tooling for agentic proving systems \cite{minif2f-v1, minif2f-v2, tsoukalas2024putnambenchevaluatingneuraltheoremprovers, requena2026minimalagentautomatedtheorem, liu2026numinaleanagentopengeneralagentic}. Our work supports this direction by aligning informal and formal statements with their dependency context.

% \subsubsection{Lean Graph}
% Dependency extraction for formal libraries is well studied \cite{alama-formal-deps}, and several Lean tools build declaration-level graphs. Jixia, the extractor behind LeanSearch~v2, emits a raw post-elaboration constant-reference graph; Lean~Atlas splits type- and value-level edges and adds confidence scores, \texttt{sorry}-status, reachability, and visualization; and \texttt{leanblueprint} links LaTeX exposition to Lean through hand-annotated \texttt{\textbackslash lean\{\}} references \cite{jixia, yanahama2026leanatlas, massot2020leanblueprint}. LeanGraph is complementary: it classifies dependencies through Lean's kernel APIs, filters kernel artifacts with a documentation-aligned node set, and refines prior edge taxonomies into \texttt{sig}/\texttt{extends}/\texttt{field} and \texttt{proof}/\texttt{def} edges at multi-library scale.

\subsubsection{Lean Graph}
\label{sec:leangraph}
Dependency extraction for formal libraries is well studied \cite{alama-formal-deps}, and several Lean tools build declaration-level graphs. Jixia, the extractor behind LeanSearch~v2 \cite{jixia}, emits a raw post-elaboration constant-reference graph; Lean~Atlas \cite{yanahama2026leanatlas} splits type- and value-level edges and adds confidence scores, \texttt{sorry}-status, reachability, and visualization; and \texttt{leanblueprint} \cite{massot2020leanblueprint} provides author-annotated \texttt{\textbackslash lean\{\}} links between informal LaTeX statements and their formal Lean declarations. LeanGraph is complementary: it classifies dependencies through Lean's kernel APIs, filters kernel artifacts with a documentation-aligned node set, and refines prior edge taxonomies into \texttt{sig}/\texttt{extends}/\texttt{field} and \texttt{proof}/\texttt{def} edges at multi-library scale.